\documentclass[11pt,a4paper]{article}
\usepackage[margin=2.6cm]{geometry}
\usepackage{amsmath,amssymb}
\usepackage{booktabs}
\usepackage{microtype}
\usepackage{parskip}
\usepackage[colorlinks=true,linkcolor=blue!50!black,urlcolor=blue!50!black]{hyperref}
\newcommand{\code}[1]{\texttt{\small #1}}

\title{What Claude Did}
\author{Robin Langer}
\date{}

\begin{document}
\maketitle
\vspace{-2em}

We let Claude play the game and this is what he did.

\section{It played once, naively}

Fight one it did the obvious thing and hunted for a period.

\begin{center}
\begin{tabular}{ll}
\toprule
$01{:}09$ & ``Exact period $5$ locked in from index $4$.'' \\
$01{:}51$ & ``Period $12$ and period $17$ disagree on the very next beat. \\
          & \ That's a free experiment --- block and look.'' \\
$02{:}10$ & ``Confirmed: period $12$ from index $9$.'' \\
$02{:}26$ & ``Fight 1 lost.'' \\
\bottomrule
\end{tabular}
\end{center}

Each period fitted everything seen. Period $12$ was even put through a test
designed to separate it from period $17$, passed, and then survived twenty-eight
further beats before it broke. The fight was lost at beat $60$.

\section{The tell}

Then it called \code{reveal}, looked at the whole schedule, and wrote this:

\begin{quote}
The tape is a \textbf{mechanical (Sturmian) word}: beat $i$ is OPEN iff
$\mathrm{frac}(i \cdot a + b) < a$. A single $(a,b)$ pair reproduced all $120$
beats with zero errors. That's why every ``period'' I found held for a while
then broke --- irrational rotations produce long \emph{almost}-periodic
stretches (the continued-fraction convergents $5/12$, $17/41 \ldots$) that
eventually slip.
\end{quote}

Nothing in its brief had mentioned Sturmian sequences, or irrational slopes, or
that the schedules were generated at all. It was given the rules of a fight and
four commands.

The second sentence is the fit. The third is the insight, and it is a different
kind of thing: \textbf{it diagnosed the generator from the pattern of its own
failures}, not from the data.

Look at what it had failed with. Periods $5$, then $12$, then $17$. That
monster's slope was $\alpha \approx 0.4142$, and
\[
  \sqrt{2} - 1 = 0.41421\ldots = [0;2,2,2,\ldots],
  \qquad
  \text{convergents } \tfrac12,\ \tfrac25,\ \tfrac{5}{12},\ \tfrac{12}{29},\ \ldots
\]
Its first two hypotheses are the denominators of successive convergents, and
$17/41$ is the mediant of $5/12$ and $12/29$ --- the next best rational
approximation between them. It had not been finding wrong periods. It had been
computing, one at a time and without realising, the continued-fraction expansion
of an irrational number.

A rational schedule has one period and keeps it. A sequence of periods that each
hold and then slip, with denominators climbing $5, 12, 17, 41$, is what an
irrational rotation looks like from inside. That is the signature, and it is
visible only in the failures.

\section{What followed}

Having named the class it stopped playing and wrote a program.

For a candidate slope $a$, one observed beat confines the phase $b$ to an arc of
the circle. A history confines it to an intersection of arcs. So rather than
choose a period, \code{engine.ts} keeps \emph{the set of every $(a,b)$ still
consistent with everything seen} --- and attacks only on beats where every
surviving pair agrees the monster is open. Where they disagree it returns
\code{"?"}, and \code{"?"} means block.

That is the $01{:}51$ idea --- disagreement is a free experiment --- promoted
from a trick into a policy.

\begin{center}
\begin{tabular}{lll}
\toprule
$03{:}57$ & \code{engine.ts} & the model, $56$ lines \\
$04{:}34$ & \code{live.ts}   & an autopilot: block twelve beats, then fit and predict \\
$04{:}34$--$07{:}58$ & & \textbf{fights 2--10, played by the program. Nine wins} \\
\bottomrule
\end{tabular}
\end{center}

All nine at full health, with not one mistaken swing. For scale, on these
monsters swinging every beat wins $0\%$, swinging only on guaranteed openings
wins $0\%$, and a fixed table on the last four beats wins $47\%$.

\section{And it measured them}

It then checked the claim properly, brute-forcing $(a,b)$ against all ten
schedules --- twenty thousand slopes by four hundred phases --- and got zero
errors on $1200$ beats. It confirmed the identification three further ways, none
suggested to it: complexity exactly $m+1$ at every window length, balance one,
and gaps between open beats taking only two consecutive values. Those three
properties are not evidence for a Sturmian word; they are the definition of one.

And it read off each monster's angle:
\[
  0.385,\; 0.414,\; 0.418,\; 0.433,\; 0.434,\; 0.455,\; 0.466,\; 0.472,\;
  0.474,\; 0.476 .
\]
The generator draws its strike rate uniformly from $[0.52, 0.62]$, so the open
rate is uniform on $[0.38, 0.48]$. It recovered the interval from ten samples,
and marked that --- correctly --- as the least confident of its claims.

\end{document}
